\section{Native ampleness on the same proper Cartier cut from integral induction}

This slice transports an intermediate lower-dimensional induction conclusion.
It does not assume the desired ampleness of the cut. The final numerical
Nakai induction must itself supply the intermediate conclusion below.
All schemes and line pullbacks remain the original ones.

\begin{lemma}[Every literal integral stratum of the original cut is ample]
\label{lem:seshadri-original-proper-cut-strata-ample}
\lean{EnvelopeProperty.SeshadriAmplenessCriterionProperCutAmpleness.sectionCut_strataAmple_of_properIntegralClosedAmple}
\leanok
\uses{thm:line-section-zero-ideal-cartier,thm:fc-curve-cartier-ideal-finite,thm:scheme-point-closure-integral,lem:tc-seven-canonical-power-transport}
Let $p:X\to\operatorname{Spec}k$ be the original projective morphism over
a field, with $X$ integral. Let $L,A$ be original invertible modules and
$s:\mathcal O_X\to A$ an original nonzero section. Set $I=I(s)$ and
$g:D=I.\mathrm{subscheme}\hookrightarrow X$ for its literal original zero cut.
Assume the following intermediate integral induction conclusion: for every
original integral scheme $Y$ and every actual closed immersion $j:Y\to X$
which is not surjective, the same original $j^*L$ is natively ample over
$\operatorname{Spec}k$ through $j\circ p$.

For each original $\eta\in D$, put
\[
 J_\eta=\operatorname{vanishingIdeal}_D
     \bigl(\overline{\{\eta\}}\bigr),\qquad
 a_\eta:J_\eta.\mathrm{subscheme}\hookrightarrow D.
\]
Then the original $a_\eta^*(g^*L)$ is natively ample for the actual
structure map $a_\eta\circ g\circ p$. There is no integrality or
reducedness hypothesis on $D$ and no supplied line comparison.
\end{lemma}
\begin{proof}
The actual nonzero original section makes $I$ a Cartier ideal by the
existing theorem. The existing proper-support theorem therefore gives
$I.\mathrm{support}\ne X$. The native range formula for $g$ identifies
its range with this same support, so $g$ is not surjective.

Fix an original point $\eta$ of $D$. The existing reduced point-closure
producer proves that the literal $J_\eta.\mathrm{subscheme}$ is integral.
Its native inclusion $a_\eta$ and $g$ are closed immersions, so their
actual composite $j=a_\eta\circ g$ is a closed immersion into $X$.
If $j$ were surjective, $g$ would be surjective, contrary to the preceding
properness. Apply the intermediate integral induction conclusion to
this actual composite. It gives native ampleness of $j^*L$, on the
same original integral point closure and with its actual base triangle.

The canonical composition comparison for module pullback gives the
whole isomorphism
\[
 c_\eta:j^*L\simeq a_\eta^*(g^*L).
\]
Unpack the obtained native ample datum: a positive integer $r$, its
actual closed projective embedding $q$, its same original base triangle,
and its whole $\mathcal O(1)$ comparison with $(j^*L)^r$.
Compose this comparison with the existing actual tensor-power
isomorphism $c_\eta^{\otimes r}$. Invertibility of
$a_\eta^*(g^*L)$ follows by the two actual pullbacks of the original
invertibility of $L$. Associativity of the original scheme maps rewrites
the base triangle to $a_\eta\circ(g\circ p)$. These constructed data
are exactly native ampleness of the stated literal line. No desired
ampleness or matching certificate has been inserted.
\end{proof}

\begin{theorem}[The same original line is ample on the whole original cut]
\label{thm:seshadri-original-proper-cut-native-ampleness}
\lean{EnvelopeProperty.SeshadriAmplenessCriterionProperCutAmpleness.isAmple_sectionCut_of_properIntegralClosedAmple}
\leanok
\uses{lem:seshadri-original-proper-cut-strata-ample,thm:seshadri-original-strata-native-ampleness}
With exactly the preceding original data and intermediate integral
induction conclusion, the same original $g^*L$ is natively ample on the
whole original $D$ over $k$. In particular $D$ may be reducible and
nonreduced. The intermediate conclusion is not advertised as an
additional assumption of the final numerical Nakai theorem.
\end{theorem}
\begin{proof}
The native cut inclusion is closed and $p$ is projective, so the actual
composite $g\circ p$ is projective. The original $g^*L$ is invertible.
The preceding lemma supplies genuine native ampleness on every literal
reduced point closure of this same $D$, with exactly the pullbacks and
structure maps required by the already proved general-ambient strata
ampleness theorem. Apply that theorem to $g\circ p$ and the same
$g^*L$. It internally assembles the original line's positive very ample
power, giving the required native ample conclusion without replacing
$D$ by its reduction or an integral component.
\end{proof}
